\appendix
\section{Robust Sylvester projectors for symmetric $3\times 3$ tensors}
\label{sec:projector-robustness}

Let $\mathbf{X}\in\mathbb{R}^{3\times 3}$ be a symmetric second-order tensor with (real) eigenvalues
$\{\lambda_1,\lambda_2,\lambda_3\}$, assumed available from an invariant-based closed-form procedure.
The spectral decomposition reads
\begin{equation}
\mathbf{X}=\sum_{i=1}^{3}\lambda_i\,\mathbf{P}_i,\qquad 
\mathbf{P}_i=\mathbf{n}_i\otimes \mathbf{n}_i,\qquad
\sum_{i=1}^{3}\mathbf{P}_i=\mathbf{I},
\end{equation}
where $\mathbf{P}_i$ are orthogonal projectors. In anisotropic settings, the projectors can be obtained
without explicit eigenvectors by Sylvester's (Lagrange) interpolation.

\paragraph{Anisotropic (three distinct eigenvalues).}
For mutually distinct $\lambda_i$, the Sylvester projectors are
\begin{equation}
\label{eq:sylvester_Pi}
\mathbf{P}_i^{\mathrm{raw}}
=
\frac{\bigl(\mathbf{X}-\lambda_j\mathbf{I}\bigr)\bigl(\mathbf{X}-\lambda_k\mathbf{I}\bigr)}
{(\lambda_i-\lambda_j)(\lambda_i-\lambda_k)},
\qquad
(i,j,k)\ \text{a permutation of}\ (1,2,3).
\end{equation}
In finite precision arithmetic, near-degeneracy renders the denominators ill-conditioned. We therefore
regularize the denominators in a scale-aware manner.
Define the spherical/deviatoric split
$q=\tfrac13\mathrm{tr}(\mathbf{X})$ and $\mathbf{B}=\mathbf{X}-q\mathbf{I}$, and a tensor scale
\begin{equation}
\label{eq:scale_def}
s^2=\max\!\Bigl(q^2+\tfrac16\mathrm{tr}(\mathbf{B}^2),\,1\Bigr).
\end{equation}
For a scalar denominator $d$, we use the stabilized form
\begin{equation}
\label{eq:denom_stab}
d_{\mathrm{safe}}=\mathrm{sign}(d)\,\max\!\bigl(|d|,\varepsilon_d\,s^2\bigr),
\end{equation}
and replace $(\lambda_i-\lambda_j)(\lambda_i-\lambda_k)$ in~\eqref{eq:sylvester_Pi} by $d_{\mathrm{safe}}$.
To enforce the partition of unity exactly, we apply the symmetric correction
\begin{equation}
\label{eq:unity_fix}
\mathbf{P}_{\Sigma}=\sum_{i=1}^{3}\mathbf{P}_i^{\mathrm{raw}},\qquad
\mathbf{P}_{\mathrm{fix}}=\tfrac13(\mathbf{I}-\mathbf{P}_{\Sigma}),\qquad
\mathbf{P}_i^{\mathrm{full}}=\mathbf{P}_i^{\mathrm{raw}}+\mathbf{P}_{\mathrm{fix}}.
\end{equation}

\paragraph{Degeneracies: transverse and full isotropy.}
When two eigenvalues coalesce (transverse isotropy), only the projector onto the simple eigenvalue is unique,
whereas any orthonormal basis in the associated two-dimensional eigenspace is admissible. To avoid arbitrary
basis selection, we introduce transverse-isotropic (TI) projector sets by splitting the degenerate subspace
symmetrically. For example, if $\lambda_2\approx \lambda_3$ (pair $(23)$), we define
\begin{equation}
\label{eq:TI_split_23}
\mathbf{P}_1^{(23)}=\mathbf{P}_1^{\mathrm{full}},\qquad
\mathbf{P}_2^{(23)}=\mathbf{P}_3^{(23)}=\tfrac12\bigl(\mathbf{I}-\mathbf{P}_1^{\mathrm{full}}\bigr),
\end{equation}
and analogously for pairs $(12)$ and $(13)$ by selecting the corresponding unique index.
In the fully isotropic limit ($\lambda_1=\lambda_2=\lambda_3$), individual eigenprojectors are not defined;
we adopt the invariant convention
\begin{equation}
\label{eq:iso_split}
\mathbf{P}_1^{\mathrm{iso}}=\mathbf{P}_2^{\mathrm{iso}}=\mathbf{P}_3^{\mathrm{iso}}=\tfrac13\mathbf{I}.
\end{equation}

\paragraph{Smooth regime blending (robust transition).}
To obtain a numerically stable and differentiable transition between anisotropy, transverse isotropy, and
full isotropy, we use smooth ``gap'' indicators based on eigenvalue separations:
\begin{equation}
\label{eq:gap_indicator}
s_{ij}=\frac{\delta^2}{(\lambda_i-\lambda_j)^2+\delta^2},\qquad
\delta^2 = (\mathrm{tol}_{\mathrm{deg}}^2)\,s^2.
\end{equation}
Here, $s_{ij}\approx 0$ for well-separated eigenvalues and $s_{ij}\approx 1$ when $\lambda_i\approx\lambda_j$.
We define unnormalized nonnegative weights for the anisotropic (full), three TI cases, and isotropy as
\begin{align}
\tilde{w}_{\mathrm{full}}&=(1-s_{12})(1-s_{23})(1-s_{13}), \nonumber\\
\tilde{w}_{12}&=s_{12}(1-s_{23})(1-s_{13}),\quad
\tilde{w}_{23}=s_{23}(1-s_{12})(1-s_{13}),\quad
\tilde{w}_{13}=s_{13}(1-s_{12})(1-s_{23}), \nonumber\\
\tilde{w}_{\mathrm{iso}}&=s_{12}s_{23}s_{13},
\end{align}
and normalize them by $w_r=\tilde{w}_r/\sum_{r}\tilde{w}_r$ (a convex combination).
The final robust projectors are then obtained as
\begin{equation}
\label{eq:final_blend}
\mathbf{P}_i
=
w_{\mathrm{full}}\,\mathbf{P}_i^{\mathrm{full}}
+w_{12}\,\mathbf{P}_i^{(12)}
+w_{23}\,\mathbf{P}_i^{(23)}
+w_{13}\,\mathbf{P}_i^{(13)}
+w_{\mathrm{iso}}\,\mathbf{P}_i^{\mathrm{iso}}.
\end{equation}
By construction, $\sum_i \mathbf{P}_i=\mathbf{I}$ and $\mathbf{P}_i=\mathbf{P}_i^{\mathsf T}$ hold in all regimes,
and the TI and isotropic limits avoid non-uniqueness of eigenvectors while preserving rotational invariance of the
degenerate subspaces.
